
Results are presented in causal order: orbit existence, then mechanism measurement, then performance.

\subsubsection{5.1 Orbit Existence (H-E1)}

\textbf{Dataset:} seed-only variant (\texttt{dataset_mnist_seed.pt}, 976 checkpoints, Conv(8) architecture). See Section 4.2 for the architectural distinction from the hyp_rand encoder experiments.

BN-free architecture was confirmed (no running_mean / running_var keys in checkpoints). Of 500 same-accuracy-decile pairs sampled across 10 accuracy deciles, all 500 pairs showed cosine distance > 0.1:

\begin{itemize}
\item orbit_proportion = \textbf{1.000}
\item mean cosine distance = \textbf{0.768} (std = 0.033)
\item Gate threshold (cosine distance > 0.05): exceeded 15×
\end{itemize}

The permutation orbit problem is present across the full accuracy spectrum of the seed-only zoo. Permutation symmetry is a structural property shared by both the Conv(8) and Conv(32)-Conv(64) architectures; direct confirmation on the hyp_rand variant was not performed.

\subsubsection{5.2 Flat MLP Permutation Sensitivity (H-M1)}

The trained flat MLP (500,577 parameters) was evaluated on 500 stratified pairs from the hyp_rand zoo:

\begin{itemize}
\item Mean L2, equivalent pairs: \textbf{4.212}
\item Mean L2, random pairs: \textbf{6.489}
\item Sensitivity score: \textbf{0.649} (gate threshold > 0.3; exceeded 2.2×)
\item Spearman ρ on hyp_rand test set: \textbf{0.1041}
\end{itemize}

The flat MLP assigns substantially different embeddings to permutation-equivalent weight configurations. The sensitivity score of 0.649 indicates that equivalent-pair distances are 65\% of random-pair distances — the encoder does not identify permutation equivalence.

\subsubsection{5.3 NFN Permutation Sensitivity (H-M2)}

The trained NFN (521,953 parameters, best checkpoint at epoch 114) was evaluated on 500 stratified pairs:

\begin{itemize}
\item Mean L2, equivalent pairs: \textbf{2.679 × 10⁻⁸}
\item Mean L2, random pairs: \textbf{3.648 × 10⁻²}
\item Sensitivity score: \textbf{7.34 × 10⁻⁷} (gate threshold < 0.1; satisfied by factor of ~136,000)
\item Sensitivity ratio relative to flat MLP: approximately \textbf{885,000×} lower (0.6490 / 7.34 × 10⁻⁷)
\item Spearman ρ on hyp_rand test set: \textbf{0.6806} [95\% CI: 0.603, 0.748]
\end{itemize}

The near-zero sensitivity score confirms that the NFN encoder maps permutation-equivalent weight configurations to effectively identical embeddings, consistent with the equivariance guarantee of the NFN architecture.

\subsubsection{5.4 Primary Result: Symmetry Spectrum and Δρ (H-M3)}

\textbf{Disclosure on flat MLP evaluation.} The H-M3 benchmark uses an untrained flat MLP instance (random weights, 500,706 params), because the trained H-M1 checkpoint was not found at the expected path. The trained H-M1 flat MLP achieved ρ = 0.1041; the untrained instance achieves ρ = 0.1688. Both are in the range 0.10–0.17, substantially below NFN ρ = 0.6806. Accordingly, Δρ = 0.5119 should be interpreted as computed against an untrained baseline; the gap against the trained flat MLP would be slightly larger (≈ 0.576), and the gap against a well-tuned multi-layer flat MLP is not measured and could be smaller (see Limitation L2).

\textbf{Table 2.} Spearman rank correlation on MNIST-CNN hyp_rand test set (n = 338). All encoders at approximately 500K parameters.

\begin{table}[htbp]
\centering
\begin{tabular}{llll}
\toprule
Encoder & Parameters & ρ & 95\% CI \\
\midrule
Flat MLP (untrained)¹ & 500,706 & 0.1688 & [0.069, 0.273] \\
Deep Sets & 471,936 & 0.4466 & [0.344, 0.544] \\
NFN (trained, epoch 114) & 521,953 & 0.6806 & [0.603, 0.748] \\
\bottomrule
\end{tabular}
\end{table}

¹ Evaluated with random (untrained) weights; trained H-M1 instance (500,577 params) achieved ρ = 0.1041.

\textbf{Δρ(NFN − flat MLP) = 0.5119} [95\% CI: 0.381, 0.638]

The CI lower bound (0.381) exceeds the gate threshold (0.05) by approximately 7.6×. The strict ordering ρ(flat MLP) < ρ(Deep Sets) < ρ(NFN) is confirmed:

\begin{itemize}
\item ρ(flat MLP) = 0.1688
\item ρ(Deep Sets) = 0.4466 (Δ = +0.278 over flat MLP)
\item ρ(NFN) = 0.6806 (Δ = +0.234 over Deep Sets)
\end{itemize}

The Deep Sets result indicates that permutation invariance (without full equivariance) provides substantial benefit, and that equivariance provides a further, smaller increment at this capacity.

\subsubsection{5.5 Accuracy-Tier Dependence (Unexpected Finding)}

The pre-registered prediction P3 anticipated mid-tier dominance of the NFN advantage. The observed data shows the opposite pattern:

\textbf{Table 3.} NFN Spearman ρ by accuracy tier on hyp_rand test set. Bootstrap CIs not computed at tier level due to small per-tier sample size (n ≈ 112–113 per tier).

\begin{table}[htbp]
\centering
\begin{tabular}{lll}
\toprule
Accuracy Tier & NFN ρ & n \\
\midrule
Low (bottom third) & 0.856 & 113 \\
Mid (middle third) & 0.317 & 113 \\
High (top third) & −0.314 & 112 \\
\bottomrule
\end{tabular}
\end{table}

Pre-registered prediction P3 (mid-tier dominance) is not supported. The NFN achieves its highest correlation for low-accuracy models and negative correlation for high-accuracy models. Tier-level values are point estimates without bootstrap CIs; the reliability of individual tier estimates is not characterized.

